% ECONOMICS_PROBLEM: {"id": "R28-415", "title": "Scaling housing mobility programs", "jel_code": "R28", "source_id": "OP-0464"}
% !TeX program = xelatex
\documentclass[11pt,a4paper]{article}
\usepackage[margin=23mm]{geometry}
\usepackage{fontspec}
\setmainfont{Latin Modern Roman}
\usepackage{amsmath,amssymb,mathtools}
\usepackage{xcolor,enumitem,booktabs,longtable,array,xurl,hyperref}
\definecolor{muted}{HTML}{53636C}
\hypersetup{colorlinks=true,urlcolor=blue,linkcolor=blue}
\setlength{\parindent}{0pt}
\setlength{\parskip}{5pt}
\newcommand{\R}{\mathbb R}
\newcommand{\E}{\mathbb E}
\newcommand{\N}{\mathbb N}
\newcommand{\ind}{\mathbf 1}
\newcommand{\Var}{\operatorname{Var}}
\newcommand{\Cov}{\operatorname{Cov}}
\newcommand{\supp}{\operatorname{supp}}
\DeclareMathOperator*{\argmin}{arg\,min}
\DeclareMathOperator*{\argmax}{arg\,max}
\newcommand{\entrymeta}[3]{{\sffamily\small\color{muted}\textbf{\texttt{#1}}\quad|\quad #2\par #3\par}}
\newcommand{\Problemref}[1]{\texttt{#1}}
\newcommand{\JELref}[2]{\texttt{#2} (JEL \texttt{#1})}
\begin{document}
\section*{Scaling housing mobility programs}
\textbf{Problem ID:} \texttt{R28-415}\quad \textbf{JEL:} \texttt{R28}\par
% BEGIN ECONOMICS STATEMENT
\label{problem:OP-0464}
{\sffamily\small\color{muted}Original subject group: Urban, housing and spatial economics\par}
\label{definitions:problem:30007746}
\textbf{Audit classification:} Published research agenda. Removed double counting of children earnings within family monetary welfare and specified market saturation.
\subsection*{Definitions and precise research target}
\textit{Editorial formalization.} Consider eligible and noneligible baseline households in one metropolitan market. At prespecified saturation \(s\), an explicitly randomized share of eligible families receives search and moving assistance; housing, rents, schools and neighborhoods respond. For a fixed baseline household population including relevant owners and taxpayers, define monetary benefits \(b_i\) and weights \(\omega_i>0\) as follows. The proposed empirical target is \[W(s)-W(0)=E\sum_i\omega_i b_i(s).\] Identify or bound it using credible intervention variation and a declared model of price, selection and capacity responses. Financing, ownership income, resources and public services enter once. Local trial effects do not automatically identify full equilibrium responses; report selection, transport and distributional assumptions.

Saturation \(s\) is a market-level allocation probability under a fully specified lottery of offers, not an individual's treatment. Let \(\mathcal U_i^s\) be expected lifetime utility of every baseline household, eligible or otherwise, under that saturation, averaging over the lottery and equilibrium responses. Include all ownership income, child outcomes, financing and program resources in budgets or utility, using a common declared horizon/terminal value. Set \(b_i(s)\) by \(\mathcal U_i^0(b_i(s))=\mathcal U_i^s(0)\) and \(W(s)=E\sum_i\omega_i b_i(s)\), assuming unique integrable benefits. Do not add children's earnings to a family money metric that already includes the resulting consumption. Fixed saturation points are estimable under randomized market saturation and a stated between-market interference restriction; individual offers in a single low-saturation trial identify a different contrast. All expectations use a stated baseline population and probability law. Identification means every admissible data-generating model with the same observed-data law yields the same displayed target; otherwise report the identified set. Exact equations, horizons and assumptions are editorial, rather than source-given canonical conjectures.
\subsection*{Variants and assumption changes}
\begin{enumerate}
\item \textbf{Editorial capacity variant.} Compare fixed housing/school supply with endogenous supply under the same saturation lottery. Congestion and rents change \(W(s)\).
\item \textbf{Editorial distributional variant.} Change announced welfare weights and report eligible, noneligible and owner benefits separately. Positive aggregate benefits need not benefit every group.
\end{enumerate}
\subsection*{References and source audit}
\textbf{1.} Explicit agenda on mechanism weights, scaling mobility programs, and incumbent residents of revitalized areas. Locator: Journal of Economic Perspectives 35(4), 197–222; Discussion and Conclusion, pp.216–217. Full text or primary publication page inspected. URL: \url{https://lkatz.scholars.harvard.edu/sites/g/files/omnuum5961/files/lkatz/files/chyn_katz_jep_2021_all.pdf}.

\textbf{2.} Opportunity-market review exists; older review supplies explicit agenda, not this abstract. Locator: NBER Working Paper 35419, July 2026; abstract. Primary indexed metadata/abstract inspected; direct record failed. URL: \url{https://www.nber.org/papers/w35419}.
\begin{enumerate}\raggedright
\item\hypertarget{definitions:source:30007746:1}{} Eric Chyn; Lawrence F. Katz. 2021. \emph{Neighborhoods Matter: Assessing the Evidence for Place Effects}. Journal of Economic Perspectives 35(4), 197–222; Discussion and Conclusion, pp.216–217.
\url{https://lkatz.scholars.harvard.edu/sites/g/files/omnuum5961/files/lkatz/files/chyn_katz_jep_2021_all.pdf}.
DOI: \url{https://doi.org/10.1257/jep.35.4.197}.
\item\hypertarget{definitions:source:30007746:2}{} Lawrence F. Katz. 2026. \emph{Neighborhood Effects and Missing Markets for Opportunity}. NBER Working Paper 35419, July 2026; abstract.
\url{https://www.nber.org/papers/w35419}.
DOI: \url{https://doi.org/10.3386/w35419}.
\end{enumerate}

\subsection*{Common conventions: Source mathematical conventions}

These conventions apply to each entry unless explicitly replaced.

\textbf{Models and identification.} A primitive model \(m\) specifies all probability laws, feasible choices, information, equilibrium and measurement rules used by its target. The admissible class \(\mathcal M\) is fixed independently of outcomes. If \(P_O(m)\) is its observable probability law and \(\tau(m)\) the target, its identified set at observable law \(P\) is
\[
 \mathcal I_\tau(P)=\{\tau(m):m\in\mathcal M,\ P_O(m)=P\}.
\]
Point identification means that this set is a singleton. An empty set signals incompatibility of data and assumptions. Coordinate bounds need not describe the joint set. Bounds are sharp only if every retained target is attainable or arbitrarily approachable by compatible models. Finite bounds require explicit restrictions on unobserved quantities; unrestricted confounding, hidden wealth, extrapolation or arbitrary equilibrium selection can yield uninformative sets.

\textbf{Causal interventions.} A potential outcome \(Y_i(p)\) is the outcome under a complete policy regime \(p\), including timing, financing, information and market spillovers. The target population and units are fixed before treatment. With interference, use the complete assignment vector or a justified exposure mapping; individual no-interference notation alone cannot identify an economy-wide intervention. A difference of mean potential outcomes does not require the cross-world joint distribution. Conditioning on post-treatment migration, survival, enrollment or employment changes the estimand and needs separate justification.

\textbf{Statistical statements.} An estimator, confidence set, forecast or test specifies a sampling law, model class and loss/coverage criterion. Uniform \((1-\alpha)\)-coverage means \(\inf_{m\in\mathcal M}\Pr_m\{\tau(m)\in C_n\}\geq1-\alpha\), or an explicitly asymptotic version. The whole parameter space trivially has coverage, so informativeness requires a stated width, risk or power objective. A forecasting improvement is relative to a named benchmark, information set and evaluation population.

\textbf{Equilibrium and optimization.} An equilibrium comprises feasible plans, optimality, beliefs where needed, budget constraints and all market/resource clearing. The primitives or a stated benchmark define each correspondence \(\mathcal E(p;m)\). Nonemptiness is assumed or proved, never inferred just from compact policy sets. A selected-equilibrium target uses a declared selection rule; a robust target quantifies over all permitted equilibria. Use suprema or infima unless attainment is established. An \(\varepsilon\)-optimal policy is feasible and within \(\varepsilon\) of the finite optimum value. Report an empty feasible set separately.

\textbf{Welfare and accounting.} Normative utility, interpersonal normalization, social weights, numeraire, discounting and the use of public receipts are primitives. Behavioral utility need not coincide with the welfare criterion. Household welfare includes ownership income and rents once; adding profits again to compensating variation generally double-counts them. A monetary equivalent is defined by an explicit transfer and price convention, with monotonicity and range conditions. Average effects, mechanism contributions, transfers and welfare losses are different targets. Infinite sums require convergence and dynamic feasible plans require terminal or no-Ponzi conditions.

\textbf{Source versions.} A working-paper date, revision date and journal publication year are distinguished. A DOI for a journal version is not automatically the DOI of an earlier manuscript. Locators refer to the stated version and printed pages where specified. A metadata-only check does not verify a claim about a full-text conclusion. Every entry's source subsection narrows the provenance claim accordingly.
% END ECONOMICS STATEMENT
\end{document}
